% entries/parameters-vs-architecture.tex
\vocabentry{Parameters versus architecture (the quiver)}
{The architecture is the shape of a network: the number of layers, the width $n_\ell$ of each, which entries of each $W^{(\ell)}$ are free, tied together, or fixed at zero, and the nonlinearity. The parameters $\theta$ are the numerical values of the free entries. Training changes $\theta$ and never the architecture.}
{Following Armenta and Jodoin (2021), an architecture is a quiver $Q$ --- a directed graph with one vertex per neuron and one arrow per weight --- with an activation attached to each hidden vertex, and a parameter vector is a representation of $Q$: a real number on every arrow and a bias on every vertex; the width is the dimension vector. The quiver isomorphisms fixing the input and output vertices --- permutations of the hidden vertices of a layer and, for ReLU, the positive rescalings $\lambda_v$ at a hidden vertex $v$ (incoming arrows times $\lambda_v$, outgoing arrows times $\lambda_v^{-1}$) --- act on representations without changing $f_\theta$; HAaNE~I calls the largest group that fixes $f_\theta$ the function-stabiliser. Two architectures have non-isomorphic quivers, so no group relates their parameters; the knowledge matrix compares them anyway because its shape depends only on $(C,d)$.}
{The arms of HAaNE~II are architectures (CNN versus MLP; ReLU, GELU, tanh; standard, NTK, $\mu$P parameterisation) crossed with a width and learning-rate grid, \texttt{slurm/configs/width\_lr*.sh}; HAaNE~I compares three ImageNet architectures.}
{\texttt{tests/test\_od4\_invariance.py}: invariance of $\KM$ under hidden-neuron permutation and positive rescaling, non-invariance under rotations; HAaNE~I Study~1 (permutation, teleportation).}
